\documentclass[10pt]{article}
\usepackage[margin=0.78in]{geometry}
\usepackage{amsmath,amssymb,amsthm,mathtools}
\usepackage{graphicx}
\usepackage{xcolor}
\usepackage{microtype}
\usepackage[hidelinks]{hyperref}
\usepackage{enumitem}
\setlist{nosep,leftmargin=*}
\setlength{\parindent}{0pt}
\setlength{\parskip}{5pt}
\newtheorem{theorem}{Theorem}
\newtheorem{lemma}[theorem]{Lemma}
\newtheorem{corollary}[theorem]{Corollary}
\newcommand{\T}{\mathbb T}
\newcommand{\C}{\mathbb C}
\newcommand{\PP}{\mathbb P}
\newcommand{\A}{\mathcal A}
\newcommand{\ran}{\operatorname{ran}}
\newcommand{\diag}{\operatorname{diag}}

\title{The Complex Almost-Injectivity Threshold Is Exactly \(2d\)\\
\large A degree obstruction for every critical rank-one frame}
\author{Automated proof attempt for arXiv:1403.1458\\
\texttt{agent\_lane\_03} \quad model: GPT5.6}
\date{11 August 2026}

\begin{document}
\maketitle

\begin{center}
\fcolorbox{blue!45!black}{blue!4}{\parbox{0.92\linewidth}{
\textbf{Status: full result, pending human verification.}
For every standard complex measurement system with \(N\le 2d-1\), almost
every signal has a nontrivial intensity twin.  Combining this universal
necessity with the known generic sufficiency at \(N\ge 2d\) proves Conjecture
2 of the source paper.}}
\end{center}

\section*{Source question}

Mixon defines almost injectivity as uniqueness modulo global phase on an open
dense set and asks for a phase transition at \(N=2M\).  The excerpt is from
printed pages 22--23, Section 3.2 of the source paper.

\begin{center}
\includegraphics[width=0.93\linewidth]{question_crop.png}
\end{center}

\section*{Result}

For \(\Phi=(\varphi_j)_{j=1}^N\subset\C^d\), write
\[
 \A_\Phi(x)=\bigl(|\langle x,\varphi_j\rangle|^2\bigr)_{j=1}^N,
 \qquad x\sim y\iff y=\omega x\quad(\omega\in\T).
\]

\begin{theorem}[Universal failure at and below the critical count]
\label{thm:main}
Let \(d\ge2\) and \(N\le2d-1\).  For every
\(\Phi\subset\C^d\), there is an open dense conull set
\(G\subset\C^d\) such that every \(x\in G\) has a \(y\not\sim x\) with
\(\A_\Phi(y)=\A_\Phi(x)\).  Consequently \(\A_\Phi\) is neither almost
injective in the source's open-dense sense nor phase retrievable almost
everywhere.
\end{theorem}

This is stronger than merely producing one bad signal: at the proposed lower
threshold, ambiguity is the generic and almost-everywhere behavior for
\emph{every} frame.

\section*{Proof intuition}

Normalize a spanning frame to a Parseval frame.  Its analysis map identifies
signals with a \(d\)-plane \(L\subset\C^{2d-1}\), and the measurements become
the coordinate magnitudes of points of \(L\).  Fix a point \(z\in L\).
Changing its coordinate phases traces a torus of real dimension \(2d-2\).
Membership back in \(L\) imposes \(d-1\) complex equations, also \(2d-2\)
real equations.  Thus the critical ambiguity problem is the zero set of a
square map
\[
 \T^{2d-2}\longrightarrow\C^{d-1}\cong\mathbb R^{2d-2}.
\]
The original point is always a zero.  If that zero is nondegenerate, it cannot
be the only one: a map from a closed manifold into Euclidean space has degree
zero, while one regular preimage would have mod-2 degree one.  If the zero is
not generically nondegenerate, the normalized intensity map has deficient
maximal rank, and the constant-rank theorem gives positive-dimensional
fibers.  Either way, generic uniqueness is impossible.

\section*{Proof}

\subsection*{1. Parseval and projective reductions}

If the measurement vectors do not span \(\C^d\), a nonzero common-kernel
vector may be added to every signal outside a proper subspace without changing
any measurement.  Hence assume they span.  Let
\(S=\sum_j\varphi_j\varphi_j^*\), put
\(g_j=S^{-1/2}\varphi_j\), and replace \(x\) by \(S^{1/2}x\).  Then
\[
 |\langle S^{1/2}x,g_j\rangle|=|\langle x,\varphi_j\rangle|,
 \qquad \sum_j g_jg_j^*=I_d.
\]
Thus no fiber changes under an invertible signal coordinate change.  The
analysis operator
\[
 V:\C^d\to\C^N,\qquad Vu=(\langle u,g_j\rangle)_j
\]
is an isometry.  Set \(L=\ran V\).

On complex projective space define the normalized coordinate-intensity map
\[
 q:\PP(L)\longrightarrow\Delta^{N-1},\qquad
 q([z])=\frac{(|z_1|^2,\ldots,|z_N|^2)}{\|z\|^2},
\]
where \(\Delta^{N-1}=\{a\in\mathbb R^N:a_j\ge0,\ \sum a_j=1\}\).
It is real analytic, and \(\dim_{\mathbb R}\PP(L)=m:=2d-2\).

Let \(r\) be the maximum rank of \(dq\).  Its maximal-rank locus is open,
dense, and conull: the rank-drop locus is the zero set of the nonzero
real-analytic section \(\bigwedge^r dq\) (with the constant-map case
immediate).  If \(r<m\), then near every maximal-rank point the constant-rank
theorem gives a local fiber of dimension \(m-r>0\).  This proves the theorem
whenever \(N-1<m\), hence whenever \(N<2d-1\), and also handles the singular
alternative at \(N=2d-1\).

It remains to treat \(N=2d-1\) and \(r=m=N-1\).

\subsection*{2. The square phase map}

The full-rank locus of \(q\) is open dense and conull.  No coordinate
functional can vanish identically on \(L\), since then \(q\) would land in a
face of dimension \(N-2<m\).  Hence the projective points whose coordinates
are all nonzero also form an open dense conull set.  Fix a unit representative
\(z\) in the intersection of these two sets.

Choose an isometry \(K:\C^{d-1}\to\C^N\) with
\(\ran K=L^\perp\).  Fixing the first coordinate phase removes global phase,
and define
\[
 \Psi_z:\T^{N-1}\longrightarrow\C^{d-1},\qquad
 \Psi_z(\lambda_2,\ldots,\lambda_N)
 =K^*\diag(1,\lambda_2,\ldots,\lambda_N)z.
\]
The identity \({\bf1}\) is a zero because \(z\in L=\ker K^*\).  Any other
zero yields another vector of \(L\) with exactly the same coordinate
magnitudes.

\begin{lemma}[Transversality equivalence]
\label{lem:transverse}
The differential \(dq_{[z]}\) is nonsingular if and only if the zero
\({\bf1}\) of \(\Psi_z\) is nondegenerate.
\end{lemma}

\begin{proof}
Represent a tangent class at \([z]\) by \(\delta\in L\) orthogonal to \(z\).
Then
\[
 (dq_{[z]}\delta)_j=2\operatorname{Re}(\overline z_j\delta_j).
\]
Since every \(z_j\ne0\), a kernel vector has the form
\(\delta_j=i h_jz_j\) with \(h_j\in\mathbb R\).  Subtracting the global
infinitesimal phase \(ih_1z\) gives the same projective tangent class and
sets \(h_1=0\).  The remaining condition \(\delta\in L\) is exactly
\[
 K^*i\diag(0,h_2,\ldots,h_N)z=0,
\]
which is the kernel equation for \(D\Psi_z({\bf1})\).  Conversely, a vector
in this latter kernel gives a zero tangent class for \(dq\); because
\(h_1=0\) and every coordinate of \(z\) is nonzero, the class is trivial only
when all \(h_j=0\).  Both real vector spaces have dimension \(N-1=2d-2\),
so injectivity is equivalent to nonsingularity.
\end{proof}

\subsection*{3. A Euclidean-target degree obstruction}

\begin{lemma}[One regular zero is impossible]
\label{lem:degree}
Let \(n>0\) and let \(F:\T^n\to\mathbb R^n\) be smooth.  A nondegenerate
zero of \(F\) cannot be its unique zero.
\end{lemma}

\begin{proof}
Suppose \(x_0\) were the unique zero and \(DF(x_0)\) were invertible.  Compose
\(F\) with the one-point compactification
\(\mathbb R^n\hookrightarrow S^n\).  The resulting map misses infinity and
is therefore null-homotopic, so its mod-2 degree is zero.  But zero is a
regular value with the single preimage \(x_0\), so its mod-2 degree is one, a
contradiction.
\end{proof}

At our chosen \([z]\), Lemma~\ref{lem:transverse} makes \({\bf1}\) a
nondegenerate zero of \(\Psi_z\).  Lemma~\ref{lem:degree} supplies a second
zero \(\lambda\).  Put
\(w=\diag(1,\lambda_2,\ldots,\lambda_N)z\).  Then
\(K^*w=0\), so \(w\in L\), and \(|w_j|=|z_j|\) for all \(j\).  If
\(w=\omega z\), the first nonzero coordinate forces \(\omega=1\), and the
remaining nonzero coordinates force \(\lambda={\bf1}\), contrary to the
choice of the second zero.  Thus \([w]\ne[z]\).

We have proved ambiguity throughout an open dense conull subset of
\(\PP(L)\).  Lifting projective points through radius and global phase, then
undoing the invertible Parseval coordinate change, gives the asserted open
dense conull set in \(\C^d\).  This proves Theorem~\ref{thm:main}. \qed

\section*{The phase transition}

\begin{corollary}[Conjecture 2 of arXiv:1403.1458]
Standard complex almost injectivity exhibits a phase transition at
\(N=2d\).
\end{corollary}

\begin{proof}
Theorem~\ref{thm:main} proves part (a): no system below \(2d\) is almost
injective.  Balan--Casazza--Edidin prove that for every \(N\ge2d\), a generic
complex frame has a dense open set of uniquely recoverable rays
\cite[Theorem 3.4]{BCE}.  This is the part (b) already cited as known in the
source paper \cite[pp.~22--23]{Mixon}.
\end{proof}

\section*{Checks, scope, and literature status}

The proof is topological and does not depend on computation.  The included
script \path{code/verify_torus_obstruction.py} deterministically builds
random critical Parseval frames in dimensions \(2,3,4\), checks both
Jacobian ranks in Lemma~\ref{lem:transverse}, and numerically finds multiple
phase zeros.  It is a stress test only.

The rank-one hypothesis is essential.  Huang--Rong--Wang--Xu construct
\(2d-1\) \emph{general} Hermitian quadratic measurements with PR-ae, while
explicitly recording the exceptional \(2d-1\) standard rank-one problem as
open and proving only generic failure \cite[Section 5, Theorem 5.3]{HRWX}.
Their construction does not have the plane--coordinate-torus structure used
above.

For novelty triage, exact-phrase and citation searches were run through
11 August 2026 for ``\(2d-1\)'', ``PR-ae'', ``standard phase retrieval'', and
``almost injective'', including the citing literature of \cite{HRWX}.  No
later resolution of the exceptional rank-one case was found.  This bounded
search supports human review; it is not an exhaustive priority claim.

\begin{thebibliography}{9}
\bibitem{Mixon}
D. G. Mixon, \emph{Phase Transitions in Phase Retrieval}, arXiv:1403.1458
(2014), especially Section 3.2 and Conjecture 2.  Included as
\texttt{source\_paper.pdf}.

\bibitem{BCE}
R. Balan, P. Casazza, and D. Edidin, \emph{On signal reconstruction without
phase}, Appl. Comput. Harmon. Anal. 20 (2006), 345--356; arXiv:math/0412411.
Included under \texttt{references/}.

\bibitem{HRWX}
M. Huang, Y. Rong, Y. Wang, and Z. Xu, \emph{Almost Everywhere Generalized
Phase Retrieval}, Appl. Comput. Harmon. Anal. 50 (2021), 16--33;
arXiv:1909.08874.  Included under \texttt{references/}.
\end{thebibliography}

\end{document}
